\documentclass{article}
\usepackage{amsmath}
\usepackage[margin=1in]{geometry}
\title{Grip Hand Surface Penetration Diagnostic}
\author{UpstreamDrift: Appearance-Only Research Reference}
\date{2026-11-19}
\begin{document}
\maketitle
\section{Scope, Frame, and Assumptions}
This editable standalone reference accompanies PR 15142 and issue 15141;
it is separate from the canonical engineering design manual.
Input coordinates are metres in the hand-grip frame: the origin is the hand
centre on the shaft, $z$ points toward the club head, and $x,y$ span the
grip cross-section. Appearance meshes are not physical contact capsules.
Input target surfaces must have finite vertices, nondegenerate triangles,
and exactly two oppositely directed incident faces per edge. Global winding
reversal is allowed. The diagnostic assumes a non-self-intersecting boundary;
the edge checks alone do not establish that assumption.
\section{Containment and Distance}
For query $q$ and target triangle vectors $a,b,c$ measured from $q$, let
\[
 \omega=2\operatorname{atan2}\left(a\cdot(b\times c),
 |a||b||c|+(a\cdot b)|c|+(b\cdot c)|a|+(c\cdot a)|b|\right).
\]
The point is inside when $|\sum\omega|>2\pi$. Its unsigned distance $d$
is the minimum over all triangle faces and edges: use orthogonal plane
projection when it lies inside a triangle, otherwise the closest point on
each clamped edge segment. Signed distance is $-d$ inside, $+d$ outside.
Face contact has zero distance. Work is one query point against all target
triangles at a time, with no convex-hull substitution or vertex-distance proxy.
For a source vertex set $V$, the reported fraction is
\[
 f=|V|^{-1}\sum_{q\in V}[s(q)<-0.002\;\mathrm{m}].
\]
For the two thumb surfaces the minimum signed distance implements containment
in either surface. The configured fraction limits remain 0.02 for the index
versus thumb and 0.11 for proximal phalanges versus palm or dorsum.
The depth threshold is explicit; shallow intersections are not claimed absent.
\section{Failure and Correction}
At source head \texttt{7e06e287f5}, the vertex-proximity guard passes despite
lead middle proximal vertices reaching 6.271 mm into the dorsum. The new
signed-distance regression fails on both hands: maximum deep fraction is
0.147619, above 0.11. Analytic cube tests distinguish a deeply interior point,
outside points, edge distance, and face contact. A concave L prism verifies
that a notch inside the convex hull is still outside the actual surface.
Reversed global winding preserves distances; open surfaces are refused.

Moving the rounded palm and dorsum sweep ends ten degrees back from the MCP
azimuth removes their deep overlap with proximal capsules at the tested
nominal geometry. The finger generators, poses and physical capsule parameters
are unchanged. Palm summaries and both OpenSim hand asset copies are rebaked.
This is a visual shell construction adjustment, not an anthropometric fit.
\section{Reproduction and Limits}
Run the unit files \texttt{test\_mesh\_penetration.py} and
\texttt{test\_grip\_hand\_part\_penetration.py} with pytest under Python 3.12,
with the pinned Tools submodule initialised. The latter covers both hands at
12.7, 14.0 and 14.5 mm grip radii using the default overlap pose.
Run \texttt{render\_grip\_hand\_evidence.py --write-receipt} with the Agg
backend for the static views. Each additional acceptance-view canvas is
1200 square pixels, includes the grip, and uses three camera views at each
visual roll of $-20,0,+20$ degrees. These are visual roll offsets, not
independently qualified physical shaft-roll states.

Vertex sampling does not bound triangle-interior penetration or establish
intersection volume. Other grip styles, parameter extremes, native runtime
acceptance, anatomical quality, human visual approval, and independent review
remain separate gates. No physical or full-epic acceptance follows from a
passing diagnostic or an image. Compilation of this reference is not asserted.
\end{document}
